\section{Introduction}\label{introduction}

This paper gives a machine-checked account of software architectural quality in terms of degrees of freedom (DOF). All claims are verified in Lean~4 with zero \texttt{sorry} placeholders.

\subsection{Formalization Statistics}

\begin{center}
\begin{tabular}{lr}
\hline
\textbf{Metric} & \textbf{Value} \\
\hline
Lines of Lean 4 (this paper) & \LeanTotalLines{} \\
Theorems/lemmas & \LeanTotalTheorems{} \\
\texttt{sorry} placeholders & \LeanTotalSorry{} \\
Proof files & \LeanTotalFiles{} \\
\hline
\textbf{Dependencies (live imports)} & \\
\texttt{AbstractClassSystem} & \LeanDepLinesAbstractClassSystem{} lines, \LeanDepTheoremsAbstractClassSystem{} theorems \\
\texttt{Ssot} & \LeanDepLinesSsot{} lines, \LeanDepTheoremsSsot{} theorems \\
\texttt{DecisionQuotient} & \LeanDepLinesDecisionQuotient{} lines, \LeanDepTheoremsDecisionQuotient{} theorems \\
\hline
\end{tabular}
\end{center}

All proofs build with \texttt{lake build}. Axiom dependencies verified via \texttt{\#print axioms}: \texttt{propext}, \texttt{Quot.sound}, \texttt{Classical.choice}.

\subsection{Central Result}

\begin{theorem}[Five-Way Equivalence]\label{thm:five-way}
For any architecture $A$ with $\mathrm{Cap}(A) > 0$, the following are equivalent:
\[
\underbrace{\mathrm{DOF}(A) = 1}_{\text{single source}}
\;\iff\;
\underbrace{L(A) \text{ maximal}}_{\text{engineering}}
\;\iff\;
\underbrace{\mathrm{SSOT}(A)}_{\text{epistemics}}
\;\iff\;
\underbrace{\mathrm{srank}(A) = 1}_{\text{information}}
\;\iff\;
\underbrace{\text{tractable sufficiency}}_{\text{complexity}}
\;\iff\;
\underbrace{\text{min thermo cost}}_{\text{physics}}
\]
\end{theorem}

Proved in Section~\ref{five-way-equivalence}. Machine-checked via \texttt{Leverage} $\to$ \texttt{DecisionQuotient} $\to$ \texttt{Ssot} $\to$ Mathlib.

\subsection{Contributions}

\begin{enumerate}
\item \textbf{DOF--Reliability Isomorphism (Theorem~\ref{thm:dof-reliability}):} Architecture with $n$ DOF is isomorphic to a series reliability system with $n$ components.

\item \textbf{Leverage--Error Tradeoff (Theorem~\ref{thm:leverage-error}):} $L(A_1) > L(A_2) \implies P_{\mathrm{error}}(A_1) < P_{\mathrm{error}}(A_2)$ at equal capabilities.

\item \textbf{Modification Complexity Gap (Theorem~\ref{thm:leverage-gap}):} Expected modification cost is proportional to DOF at fixed capabilities.

\item \textbf{Physical Edit-Energy Floor (Theorem~\ref{thm:physical-energy-floor}):} Minimum edit-energy $= j_\mathcal{M} \cdot \mathrm{DOF}(A)$ under Landauer calibration. No P~$\neq$~coNP assumption. \leanmeta{\LH{thermodynamic\_selection\_unconditional}, \LH{srank\_energy\_lower\_bound}}

\item \textbf{England Replication Inequality (Theorem~\ref{thm:england}):} $\Delta S_{\min}(k) - \Delta S_{\min}(1) \geq k_B \ln k$. Proof: a $k$-variable system has $2^k$ states (counting); entropy gap follows from $k \leq 2^{k-1}$ (induction). \leanmeta{\LH{england\_replication\_inequality}}

\item \textbf{Optimal Architecture (Theorem~\ref{thm:optimal}):} $f(R) = \arg\max_{A:\,\mathrm{Cap}(A) \supseteq R} L(A)$ minimizes expected error and is computable.
\end{enumerate}

\paragraph{What is new.} \texttt{AbstractClassSystem} and \texttt{Ssot} establish two of the five equivalences. This paper adds leverage maximization and proves it equivalent to the remaining three. The England inequality is new: the only physics input is the Landauer constant $k_B$; the bound reduces to $k \leq 2^{k-1}$.

\subsection{Dependency Chain}

\begin{enumerate}
\item \texttt{AbstractClassSystem}: axis orthogonality $\to$ error independence (Theorem~\ref{thm:error-independence})
\item \texttt{Ssot}: DOF~$= 1$ $\iff$ coherence $\to$ second equivalence
\item \texttt{DecisionQuotient}: tractability boundary and thermodynamic lift $\to$ fourth and fifth equivalences
\item \texttt{Leverage} (this paper): leverage maximization $\iff$ all of the above
\end{enumerate}

\subsection{Scope}

Leverage measures capability-to-DOF ratio. Performance, security, and other dimensions are orthogonal. Error independence is derived from axis orthogonality in \texttt{AbstractClassSystem}, not assumed. Physical claims use one explicit constant ($j_\mathcal{M} > 0$) and no additional axioms beyond the Landauer calibration.

\subsection{Organization}

Section~\ref{foundations} defines Architecture, DOF, Capabilities, and Leverage. Section~\ref{probability-model} derives the error model. Section~\ref{main-theorems} proves decidability and optimality. Section~\ref{instances} gives leverage instances. Section~\ref{five-way-equivalence} proves the five-way equivalence and the England inequality. Section~\ref{appendix-lean} describes the Lean mechanization.
